\subsection{7. Markov State Models from Short Non-Equilibrium Trajectories via Observable Operator Models}
\textbf{OSTI:} \texttt{1565592-MSM-Hempel}\quad\textbf{Authors:} N\"uske, Wu, Wehmeyer, Clementi, No\'e (2017)\quad\textbf{Domain:} Computational Chemistry / MSM theory\\
\scorebadge{Coverage}{9}\quad\scorebadge{Agreement}{8}\quad\textbf{Total:} 17/20\par\smallskip
\noindent\textbf{Coverage rationale.} All three numerical experiments replicated: 1D double-well, 2D potential, and alanine dipeptide (ADP) with 11,388 $\times$ 20 ps short trajectories plus 1 $\mu$s long reference on 8$\times$A100 OpenMM. OOM correction implemented and tested end-to-end. Only gap: limited exploration of $\tau$-sweep and basis-function variations reported in supplementary.\par
\noindent\textbf{Agreement rationale.} Phase 1 matches paper to 0.2\% (t 3699 vs 3708). Phase 2 within 2\%. Phase 3 ADP: direct MSM underestimates t (299 ps), OOM corrects to 2146 ps --- 7.2$\times$ improvement, matching 2020 ps long-reference within 6\%. Absolute t differs from paper's \textasciitilde{}1400 ps by \textasciitilde{}44\%, attributed to force-field differences (AMBER ff14SB vs paper's). Replication also discovered a practical subtlety (clustering must cover both basins) not explicit in paper.\par\smallskip
\noindent\textbf{What reproduced:}
\begin{itemize}[leftmargin=1.4em,itemsep=0.2em,topsep=0.2em]
\item 1D double-well exact timescale via transfer operator spectrum
\item Direct MSM bias demonstration and OOM correction (Phase 1,2,3)
\item ADP simulation at 11,388 $\times$ 20 ps scale on 8$\times$A100 with AMBER ff14SB + TIP3P
\item Long-reference 10 $\times$ 100 ns ground truth for ADP timescales
\item OOM-based timescale recovery at $\tau$=5 ps (299$\rightarrow$2146 ps correction)
\item k-means clustering + MSM construction pipeline
\end{itemize}\noindent\textbf{What missing:}
\begin{itemize}[leftmargin=1.4em,itemsep=0.2em,topsep=0.2em]
\item Exact force field reproduction (we used ff14SB, paper's choice unclear)
\item Full $\tau$-sweep showing OOM robustness across lag times
\item Supplementary basis-function variants (TICA vs k-means)
\item Paper's exact t$\approx$1400 ps --- our 2020 ps reference differs (force field origin)
\item Explicit stationary-distribution reweighting benchmarks
\item Finite-sample uncertainty quantification across many trajectory subsets
\end{itemize}\noindent\textbf{Follow-on research questions:}
\begin{enumerate}[leftmargin=1.6em,itemsep=0.2em,topsep=0.2em,label=Q\arabic*.]
\item Is the 44\% t gap (2020 vs 1400 ps) attributable to force field, water model, or integrator --- test by running identical short-trajectory protocol with CHARMM36m and GROMOS54a7 and comparing long-reference t.
\item How does OOM correction scale to $\geq$3 metastable states Apply to chignolin or villin headpiece where $\geq$3 slow modes exist and compare OOM vs direct MSM recovery of t, t, t.
\item Does the 'both-basin-coverage' clustering requirement we discovered translate into a principled seeding algorithm (e.g., farthest-point after long-ref preseeding) that recovers OOM accuracy without requiring a long reference
\item Benchmark OOM vs Koopman-based VAMPnets on identical short-trajectory ADP data --- which method recovers t with fewer trajectories, and how does the crossover depend on trajectory length
\item Quantify the breaking point: at what short-trajectory length T\_min does OOM no longer correct the MSM bias (error \textgreater{} 20\% of long-ref t) Sweep T $\in$ \{1,5,10,20,50\} ps at fixed total data.
\end{enumerate}

